\documentclass[12pt]{article}
\usepackage{amsmath, amssymb}
\usepackage{enumitem}
\usepackage{geometry}
\geometry{margin=1in}

\begin{document}

\begin{center}
\textbf{\Large Lecture Scribe} \\[6pt]
\textbf{Course:} CSE400 – Fundamentals of Probability in Computing \\
\textbf{Lecture 20:} Linear Transformations and Expectations of Multiple Random Variables \\
\textbf{Instructor:} Dhaval Patel, PhD \\
\textbf{Date:} March 19, 2025
\end{center}

\hrule
\vspace{8pt}

\section*{Introduction / Scope}

This lecture presents the following topics in sequence:

\begin{itemize}
\item Statistical representation of random vectors
\item Mean vector, correlation, and covariance matrices
\item Linear transformations
\item Mapping properties under \( Y = AX + b \)
\item Multivariate Gaussian distributions
\item Transition from 1D / 2D cases to \( N \)-dimensional joint PDFs
\item Special cases and properties
\item Implications of zero correlation and matrix factorization
\end{itemize}

All material below is strictly based on the provided lecture slides.  
No additional intuition, examples, simplifications, or inferred content has been introduced.  
Where mathematical expressions or steps were not visible or not explicitly stated in the lecture material, they have not been reconstructed or guessed.

\section*{Definitions and Notation}

\subsection*{Random Vector}

The lecture defines a vector of random variables as:

\begin{itemize}
\item An \( N \)-dimensional random vector
\item Example context specifies \( N = 5 \)
\end{itemize}

This vector represents multiple stochastic quantities in a system.

\subsection*{Mean Vector}

A mean vector is introduced as a statistical representation of typical system conditions.

(No explicit formula or symbolic expression was visible in the provided lecture material.)

\subsection*{Covariance Matrix}

The covariance matrix is presented with interpretation and definition context.

The lecture associates covariance with dependency relationships among variables.

Stated relationships include:

\begin{itemize}
\item More rain $\Rightarrow$ more congestion
\item More rain $\Rightarrow$ lower speed
\item More traffic $\Rightarrow$ higher pollution
\item Higher temperature $\Rightarrow$ pollution disperses
\end{itemize}

The covariance matrix is described as:

\begin{quote}
``Dependency map of the city''
\end{quote}

(No explicit matrix form or element-wise definition was visible in the provided material.)

\section*{Results / Theorems}

\subsection*{Joint Gaussian Distribution Assumption}

The lecture states that the system is assumed to follow a joint Gaussian distribution.

Implications stated:

\begin{itemize}
\item The system behaves in a structured probabilistic way
\item Characterized by:
\begin{itemize}
\item Mean (typical conditions)
\item Covariance (interdependencies)
\end{itemize}
\end{itemize}

(No explicit joint PDF formula was visible in the provided lecture content.)

\section*{Proofs / Derivations}

\subsection*{Linear Transformation Mapping}

A linear transformation is introduced in the form:

\[
Y = AX + b
\]

This mapping is used to define derived system quantities.

No step-by-step algebraic derivation of transformed expectation or covariance was visible in the provided slides and therefore is not included.

\section*{Worked Example (Case Study from Lecture)}

\subsection*{Smart City Traffic + Weather System}

The lecture introduces a case study involving prediction of traffic congestion using transportation and climate variables.

\subsubsection*{Step 1: Define \( N \) Random Variables}

\begin{itemize}
\item A vector of random variables is defined
\item Dimension specified: \( N = 5 \)
\end{itemize}

\subsubsection*{Step 2: Mean Vector}

\begin{itemize}
\item Introduced as representation of typical system behavior
\item No symbolic computation shown in visible material
\end{itemize}

\subsubsection*{Step 3: Covariance Matrix Interpretation}

Dependencies stated:

\begin{itemize}
\item Rainfall increase $\rightarrow$ congestion increases
\item Rainfall increase $\rightarrow$ speed decreases
\item Traffic increase $\rightarrow$ pollution increases
\item Temperature increase $\rightarrow$ pollution disperses
\end{itemize}

\subsubsection*{Step 4: Joint Gaussian Distribution}

Assumption:

\begin{itemize}
\item Entire system follows a joint Gaussian structure
\item Characterized using mean and covariance
\end{itemize}

\subsubsection*{Step 6: Linear Transformation (Congestion Index)}

The lecture defines a congestion index using a linear transformation.

This is stated explicitly as:

\[
\text{Congestion Index} = \text{Linear transformation of random vector}
\]

(No explicit coefficient matrix or vector form was visible.)

\section*{Inference Statements from Lecture}

If rainfall increases suddenly:

\begin{itemize}
\item Speed $\downarrow$
\item Traffic $\uparrow$
\item Pollution $\uparrow$
\end{itemize}

These directional effects are stated without formal derivation steps in the visible lecture content.

\section*{End of Scribe}

This lecture scribe:

\begin{itemize}
\item follows lecture order
\item preserves only visible notation and statements
\item includes no added theory or reconstructed derivations
\item avoids new examples or pedagogical simplifications
\item keeps unclear or missing mathematical content unfilled
\end{itemize}

\end{document}